\documentclass[12pt,a4paper]{article}

% ==================== 包 ====================
\usepackage[utf8]{inputenc}
\usepackage[T1]{fontenc}
\usepackage{ctex}                  % 中文支持
\usepackage{geometry}              % 页边距
\usepackage{graphicx}              % 图片
\usepackage{amsmath,amssymb}       % 数学公式
\usepackage{booktabs}              % 表格
\usepackage{hyperref}              % 超链接
\usepackage{listings}              % 代码
\usepackage{xcolor}                % 颜色
\usepackage{algorithm}             % 算法
\usepackage{algpseudocode}         % 伪代码
\usepackage{caption}               % 标题
\usepackage{subcaption}            % 子图
\usepackage{float}                 % 浮动体
\usepackage{fancyhdr}              % 页眉页脚

% ==================== 页面设置 ====================
\geometry{left=2.5cm, right=2.5cm, top=2.5cm, bottom=2.5cm}
\pagestyle{fancy}
\fancyhf{}
\fancyhead[C]{MC-agint: 基于大语言模型的 Minecraft 自主智能体}
\fancyfoot[C]{\thepage}

% ==================== 代码样式 ====================
\lstset{
    language=Java,
    basicstyle=\ttfamily\small,
    keywordstyle=\color{blue}\bfseries,
    commentstyle=\color{gray},
    stringstyle=\color{red},
    numbers=left,
    numberstyle=\tiny\color{gray},
    frame=single,
    breaklines=true,
    captionpos=b
}

% ==================== 超链接 ====================
\hypersetup{
    colorlinks=true,
    linkcolor=blue,
    citecolor=blue,
    urlcolor=blue
}

% ==================== 标题信息 ====================
\title{
    \textbf{MC-agint: 基于大语言模型的 Minecraft 自主智能体}\\[0.5em]
    \large 一种融合自然语言对话与自主行为的虚拟世界 AI 玩家架构
}
\author{
    赵冬冬\\
    \texttt{ZhaoDongdong0109@github}
}
\date{2026 年 5 月}

\begin{document}

% ==================== 封面 ====================
\maketitle
\thispagestyle{fancy}

% ==================== 摘要 ====================
\begin{abstract}
\noindent
传统 Minecraft 模组中的非玩家角色（NPC）通常依赖预设脚本或有限的状态机进行行为控制，缺乏自主性与自然交互能力。本文提出 \textbf{MC-agint}，一种基于大语言模型（LLM）的 Minecraft 自主智能体架构，使 AI 以独立玩家的身份存在于游戏世界中。该系统包含三个核心设计：（1）\textbf{自然语言对话系统}，将聊天与行为解耦，使 AI 以人类方式与玩家交流；（2）\textbf{自主行为引擎}，通过感知--思考--行动循环驱动 AI 自主探索、采集与建造；（3）\textbf{三层记忆架构}，采用工作记忆、长期记忆与情景记忆的分层设计，配合 LLM 驱动的记忆蒸馏机制，实现知识的持久化与智能化管理。实验表明，MC-agint 能够在玩家交互中展现出自然、连贯且具有个性的行为模式，为虚拟世界中具身 AI 智能体的研究提供了可落地的工程方案。

\vspace{0.5em}
\noindent\textbf{关键词：}大语言模型；Minecraft；自主智能体；自然语言处理；记忆系统；游戏 AI
\end{abstract}

\newpage
\tableofcontents
\newpage

% ==================== 1. 引言 ====================
\section{引言}

\subsection{研究背景}

Minecraft 作为全球最畅销的沙盒游戏之一，拥有超过 3 亿份销量，其开放世界的特性为人工智能研究提供了理想的实验平台\cite{minecraft2023}。与传统游戏不同，Minecraft 没有固定的游戏目标，玩家可以在无限生成的三维世界中自由探索、采集资源、建造建筑、与生物互动。这种开放性使得 Minecraft 成为测试具身智能体（Embodied Agent）能力的理想环境。

近年来，大语言模型（Large Language Model, LLM）的快速发展为自然语言理解和生成带来了质的飞跃\cite{brown2020language, ouyang2022training}。GPT-4、DeepSeek 等模型不仅能够进行流畅的对话，还展现出推理、规划和决策能力。这为将 LLM 应用于游戏 AI 提供了新的可能性。

\subsection{问题陈述}

现有的 Minecraft AI 方案主要存在以下局限：

\begin{enumerate}
    \item \textbf{交互不自然}：传统 NPC 依赖预设对话树或关键词匹配，无法进行自由对话。
    \item \textbf{行为机械化}：基于状态机或行为树的 AI 缺乏自主规划能力，行为模式单一。
    \item \textbf{记忆缺失}：大多数 AI 系统不具备跨会话的记忆能力，无法从历史交互中学习。
    \item \textbf{格式束缚}：要求 LLM 输出结构化 JSON 会严重限制其自然语言生成质量。
\end{enumerate}

\subsection{贡献}

本文的主要贡献如下：

\begin{enumerate}
    \item 提出一种将聊天与行为解耦的对话架构，使 AI 能够以自然语言与玩家交流，同时保持自主行动能力。
    \item 设计三层记忆系统（工作记忆、长期记忆、情景记忆），通过 LLM 驱动的记忆蒸馏实现知识的持续积累。
    \item 实现完整的 Minecraft Forge Mod，包含真实物理移动、原版配方合成、碰撞检测等游戏机制。
    \item 开源项目代码，为社区提供可扩展的 AI 智能体框架。
\end{enumerate}

% ==================== 2. 相关工作 ====================
\section{相关工作}

\subsection{游戏 AI}

游戏 AI 的发展经历了从规则系统到学习系统的演变。早期游戏 AI 主要依赖有限状态机（FSM）和行为树（Behavior Tree）\cite{millington2019artificial}，这些方法在确定性环境中表现良好，但缺乏适应性。强化学习（RL）的引入使 AI 能够通过试错学习策略\cite{mnih2015human}，Minecraft 的 Malmo 平台\cite{johnson2016malmo} 为此类研究提供了标准化接口。

\subsection{LLM 驱动的智能体}

将 LLM 作为智能体的核心控制器是近两年的研究热点。Generative Agents\cite{park2023generative} 在虚拟小镇中展示了 LLM 驱动的 NPC 如何进行可信的社会行为。Voyager\cite{wang2023voyager} 将 GPT-4 应用于 Minecraft，通过代码生成实现技能学习。AutoGPT\cite{autogpt2023} 和 BabyAGI\cite{babyagi2023} 探索了 LLM 自主任务规划的可能性。

与上述工作不同，MC-agint 的定位是\textbf{多人游戏中的 AI 玩家}，强调与真人玩家的自然交互和长期共存，而非单人任务完成。

\subsection{AI 记忆系统}

LLM 智能体的记忆管理是一个活跃的研究方向。LangChain\cite{langchain2023} 提出了短期记忆（对话历史）与长期记忆（跨会话存储）的分层架构。CoALA 框架\cite{coala2023} 将智能体记忆分为语义记忆（事实）、情景记忆（经历）和程序性记忆（技能），借鉴了认知心理学的分类方法。

MC-agint 在此基础上设计了三层记忆系统，并引入了基于 LLM 的\textbf{记忆蒸馏}机制——定期将工作记忆中的原始记录提炼为高价值的长期知识。

% ==================== 3. 系统架构 ====================
\section{系统架构}

\subsection{整体设计}

MC-agint 的系统架构如图~\ref{fig:architecture} 所示，主要由以下模块组成：

\begin{figure}[H]
    \centering
    \fbox{\parbox{0.85\textwidth}{
        \centering
        \vspace{1em}
        \textbf{MC-agint 系统架构}\\[1em]
        \begin{tabular}{|c|c|c|}
            \hline
            \textbf{对话管理} & \textbf{自主核心} & \textbf{记忆系统} \\
            \hline
            自然语言回复 & 感知--思考--行动循环 & 工作/长期/情景记忆 \\
            \hline
        \end{tabular}\\[0.5em]
        \begin{tabular}{|c|c|c|c|}
            \hline
            \textbf{动作执行} & \textbf{API 客户端} & \textbf{状态机} & \textbf{命令系统} \\
            \hline
            移动/挖掘/合成/攻击 & 异步 LLM 调用 & 行为状态管理 & /ai 管理 \\
            \hline
        \end{tabular}\\[0.5em]
        \begin{tabular}{|c|}
            \hline
            \textbf{Minecraft Forge 引擎（1.20.1）} \\
            \hline
        \end{tabular}
        \vspace{1em}
    }}
    \caption{MC-agint 系统架构图}
    \label{fig:architecture}
\end{figure}

\subsection{核心循环}

AI 的行为由一个每游戏 tick（50ms）执行一次的核心循环驱动，其伪代码如算法~\ref{alg:core_loop} 所示：

\begin{algorithm}[H]
\caption{AI 核心循环（每 Tick）}
\label{alg:core_loop}
\begin{algorithmic}[1]
\Require $tick$ 为当前游戏 tick
\State $stateMachine.tick()$
\State $memory.tick()$ \Comment{自动存盘与蒸馏检查}
\If{有异步 LLM 响应返回}
    \State $handleLLMResponse(response, sender)$
\EndIf
\If{玩家死亡}
    \State $handleDeath()$
    \State \Return
\EndIf
\State $tickPickupItems()$ \Comment{自动拾取}
\If{正在挖掘}
    \State $tickMining()$
    \State \Return
\EndIf
\If{正在吃东西}
    \State $tickEating()$
    \State \Return
\EndIf
\If{检测到危险}
    \State $handleDanger()$
    \State \Return
\EndIf
\Switch{$stateMachine.currentState$}
    \Case{IDLE}
        \If{冷却结束}
            \State $\rightarrow PLANNING$
        \EndIf
    \EndCase
    \Case{PLANNING}
        \State $prompt \gets buildPrompt(perception, memory, goal)$
        \State $LLM.callAsync(prompt)$
    \EndCase
    \Case{EXECUTING}
        \State $executeAction(action, target)$
    \EndCase
\EndSwitch
\end{algorithmic}
\end{algorithm}

% ==================== 4. 对话系统 ====================
\section{对话系统}

\subsection{聊天与行为解耦}

传统方案要求 LLM 在聊天时同时输出对话内容和行为指令（通常为 JSON 格式），这导致两个问题：

\begin{enumerate}
    \item JSON 格式约束限制了自然语言生成质量；
    \item 模型需要同时完成对话生成和行为决策两个任务，增加认知负荷。
\end{enumerate}

MC-agint 采用\textbf{解耦设计}：聊天回复使用纯自然语言，行为决策由自主思考系统独立处理。具体而言：

\begin{itemize}
    \item \textbf{聊天路径}：玩家消息 $\rightarrow$ 自然语言 prompt $\rightarrow$ AI 直接回复（纯文本）
    \item \textbf{自主路径}：冷却触发 $\rightarrow$ 行为规划 prompt $\rightarrow$ JSON 格式输出（动作 + 目标）
\end{itemize}

聊天 prompt 的设计原则是\textbf{最小约束}：

\begin{lstlisting}[caption={聊天 Prompt 核心设计},label={lst:chat_prompt}]
// 系统身份
"你是 {name}，一个生活在 Minecraft 世界里的人。
你有自己的想法、经历和感受。你不是助手，你是朋友。"

// 回复指引
"像朋友一样回复。根据对方说的话决定你的回复长短和深度——
随口打招呼就简短回应，想聊就多说几句，有感触就分享感受。
只说话，不要加任何格式标记。"
\end{lstlisting}

\subsection{目标持久化}

当 AI 在聊天中承诺执行某项任务（如"我去盖个房子"），系统会将其回复存储为\textbf{当前目标}（\texttt{currentGoal}）。在后续的自主思考中，该目标会被注入 prompt：

\begin{lstlisting}[caption={目标注入 Prompt},label={lst:goal_prompt}]
"【你之前答应要做的事】
{currentGoal}
继续做这件事，直到完成。如果已经完成或做不了，action 填 none。"
\end{lstlisting}

这种设计使 AI 能够在没有玩家持续监督的情况下完成长期任务，同时保留了玩家通过 \texttt{/ai stop} 命令中断任务的能力。

\subsection{个性化生成}

每个 AI 在首次召唤时会随机生成性格特征，由两个维度组合：

\begin{table}[H]
\centering
\caption{性格特征组合}
\label{tab:personality}
\begin{tabular}{@{}ll@{}}
\toprule
\textbf{特质} & \textbf{风格} \\
\midrule
好奇心强，喜欢探索新地方 & 说话简洁直接 \\
沉稳务实，做事有条理 & 喜欢用比喻 \\
热情开朗，喜欢交朋友 & 偶尔会冒出奇怪的想法 \\
独立自主，有自己的节奏 & 说话带点文艺范 \\
幽默风趣，经常开玩笑 & 很接地气，像朋友聊天 \\
勇敢无畏，遇到危险不退缩 & \\
细心谨慎，做事前先观察 & \\
创造力强，喜欢建东西 & \\
\bottomrule
\end{tabular}
\end{table}

% ==================== 5. 自主行为引擎 ====================
\section{自主行为引擎}

\subsection{状态机}

AI 的行为由一个有限状态机管理，包含以下状态：

\begin{table}[H]
\centering
\caption{行为状态机}
\label{tab:states}
\begin{tabular}{@{}lll@{}}
\toprule
\textbf{状态} & \textbf{描述} & \textbf{触发转换} \\
\midrule
IDLE & 空闲，等待下一次思考 & 冷却结束 $\rightarrow$ PLANNING \\
PLANNING & 调用 LLM 进行行为规划 & LLM 响应 $\rightarrow$ EXECUTING \\
EXECUTING & 执行动作 & 动作完成 $\rightarrow$ VERIFYING \\
VERIFYING & 验证动作结果 & 验证完成 $\rightarrow$ IDLE \\
CHATTING & 与玩家对话 & 对话结束 $\rightarrow$ IDLE \\
EXPLORING & 探索周围环境 & 探索结束 $\rightarrow$ IDLE \\
DANGER & 检测到危险 & 逃离 $\rightarrow$ FLEEING \\
FLEEING & 逃跑中 & 安全 $\rightarrow$ IDLE \\
\bottomrule
\end{tabular}
\end{table}

\subsection{感知系统}

AI 的感知信息（perception）包含以下维度：

\begin{itemize}
    \item \textbf{位置与维度}：当前坐标、所在维度（主世界/下界/末地）
    \item \textbf{状态}：生命值、饥饿度、时间（白天/夜晚）
    \item \textbf{环境}：脚下方块类型
    \item \textbf{实体}：15 格内的实体及其距离
    \item \textbf{背包}：物品清单与数量
\end{itemize}

\subsection{动作系统}

AI 支持的动作通过真实 Minecraft API 实现，而非简单的坐标修改：

\begin{table}[H]
\centering
\caption{动作实现方式}
\label{tab:actions}
\begin{tabular}{@{}ll@{}}
\toprule
\textbf{动作} & \textbf{实现方式} \\
\midrule
移动 & 碰撞检测 + 台阶处理 + 逐 tick 移动 \\
挖掘 & \texttt{ServerLevel.destroyBlock()}，消耗时间模拟 \\
合成 & \texttt{RecipeManager} 查询配方 + 材料消耗 + 成品给予 \\
放置 & \texttt{BlockItem.useOn()} 触发完整放置流程 \\
攻击 & \texttt{Player.attack()} + 攻击冷却 + 自动换武器 \\
吃东西 & \texttt{startUsingItem()} + 32 tick 持续时间 \\
\bottomrule
\end{tabular}
\end{table}

\subsection{危险检测}

系统在每个 tick 检测以下危险条件，触发快速响应路径（不等待 LLM）：

\begin{itemize}
    \item 生命值低于 6（3 颗心）
    \item 6 格内存在敌对生物
    \item 脚下或头部为岩浆/火焰
\end{itemize}

检测到危险后，AI 会沿威胁反方向逃跑，使用加速移动（\texttt{RUN\_SPEED}），直到脱离威胁范围。

% ==================== 6. 记忆系统 ====================
\section{记忆系统}

\subsection{三层架构}

受认知心理学启发，MC-agint 采用三层记忆架构：

\begin{figure}[H]
    \centering
    \fbox{\parbox{0.8\textwidth}{
        \centering
        \vspace{0.5em}
        \begin{tabular}{|c|c|c|}
            \hline
            \textbf{工作记忆} & \textbf{长期记忆} & \textbf{情景记忆} \\
            \hline
            最近 30 条原始思考 & 蒸馏后的重要事实 & 重要事件记录 \\
            满了自动蒸馏 & 按重要性排序 & 和谁做了什么 \\
            FIFO 淘汰 & 永久保存 & 永久保存 \\
            \hline
        \end{tabular}
        \vspace{0.5em}
    }}
    \caption{三层记忆架构}
    \label{fig:memory}
\end{figure}

\subsection{记忆蒸馏}

记忆蒸馏是 MC-agint 记忆系统的核心机制。当工作记忆达到容量上限（30 条）时，系统会异步调用 LLM 对早期记忆进行提炼：

\begin{lstlisting}[caption={蒸馏 Prompt},label={lst:distill_prompt}]
"以下是你最近的经历，请提炼出值得长期记住的事实。

最近的经历：
- {thought_1}
- {thought_2}
- ...

请输出 JSON 数组，每条包含：
- content: 事实内容（一句话）
- importance: 1-5（5=必须记住）
- type: fact / preference / knowledge / event"
\end{lstlisting}

蒸馏过程具有以下特性：

\begin{itemize}
    \item \textbf{异步执行}：不阻塞游戏主线程
    \item \textbf{保留最新}：最近 5 条记录不参与蒸馏，保证上下文连续性
    \item \textbf{去重检查}：避免重复存储相同事实
    \item \textbf{重要性排序}：长期记忆按重要性淘汰，低价值记忆被优先移除
\end{itemize}

\subsection{记忆检索}

给 prompt 注入记忆时，系统从三层中分别提取：

\begin{enumerate}
    \item \textbf{工作记忆}：最近 5 条原始思考（即时上下文）
    \item \textbf{长期记忆}：按重要性取 top 10（核心知识）
    \item \textbf{情景记忆}：最近 5 条事件（经历回顾）
\end{enumerate}

\subsection{持久化}

记忆以 JSON 格式存储于 \texttt{ai-agent-memory/} 目录，包含自动存盘机制：

\begin{itemize}
    \item 脏标记（dirty flag）：记忆变更时标记
    \item 定时保存：每 2 分钟检查脏标记并保存
    \item 关服保存：服务器关闭时强制保存
    \item 崩溃恢复：重启后自动加载，最多丢失最近 2 分钟的记忆
\end{itemize}

% ==================== 7. 实现细节 ====================
\section{实现细节}

\subsection{技术栈}

\begin{table}[H]
\centering
\caption{技术栈}
\label{tab:techstack}
\begin{tabular}{@{}ll@{}}
\toprule
\textbf{组件} & \textbf{技术} \\
\midrule
游戏版本 & Minecraft 1.20.1 \\
模组框架 & Minecraft Forge \\
开发语言 & Java 17 \\
LLM 接口 & OpenAI / DeepSeek / Ollama 兼容 API \\
构建工具 & Gradle \\
序列化 & Gson \\
\bottomrule
\end{tabular}
\end{table}

\subsection{异步 LLM 调用}

为避免 LLM 调用阻塞服务器主线程（导致游戏卡顿），所有 LLM 请求通过 \texttt{CompletableFuture} 异步执行：

\begin{lstlisting}[caption={异步调用模式},label={lst:async}]
// 发起异步调用
waitingForLLM = true;
apiClient.chatAsync(prompt).thenAcceptAsync(response -> {
    pendingLLMResponse = response;
}, serverExecutor);  // 回调在主线程执行

// 在 tick() 中检查结果
if (pendingLLMResponse != null) {
    handleLLMResponse(pendingLLMResponse, sender);
}
\end{lstlisting}

\subsection{移动系统}

AI 的移动使用 Minecraft 原生的物理系统，包括：

\begin{itemize}
    \item 碰撞检测：检查目标位置是否可通行（非固体方块）
    \item 台阶处理：自动跨上 0.6 格高的台阶
    \item 位置广播：通过 \texttt{ClientboundTeleportEntityPacket} 同步到所有客户端
    \item 朝向更新：面向移动方向
\end{itemize}

\subsection{API 容错}

LLM API 调用包含以下容错机制：

\begin{itemize}
    \item 自动重试：超时或 5xx 错误时最多重试 2 次
    \item 限流处理：429 响应时指数退避
    \item 响应解析容错：支持 JSON、Markdown 代码块、纯文本等多种格式
    \item 字段名兼容：尝试多个可能的字段名（如 \texttt{reply}、\texttt{message}、\texttt{response}、\texttt{content}）
\end{itemize}

% ==================== 8. 讨论 ====================
\section{讨论}

\subsection{优势}

\begin{enumerate}
    \item \textbf{自然交互}：解耦设计使 AI 的对话质量不受行为输出格式限制，回复更加自然流畅。
    \item \textbf{自主性}：AI 在没有玩家指令时自主探索、采集、建造，增强了世界的生动感。
    \item \textbf{持续学习}：三层记忆系统使 AI 能够积累知识，行为随时间推移更加智能。
    \item \textbf{低侵入性}：作为 Forge Mod 实现，无需修改游戏核心代码，兼容其他模组。
\end{enumerate}

\subsection{局限性}

\begin{enumerate}
    \item \textbf{LLM 依赖}：所有智能行为依赖外部 API，离线场景需要本地模型支持（如 Ollama）。
    \item \textbf{延迟}：LLM 调用的网络延迟（通常 1--3 秒）使 AI 的反应速度不如真人玩家。
    \item \textbf{成本}：频繁的 API 调用会产生费用，长时间运行的服务器成本较高。
    \item \textbf{幻觉}：LLM 可能生成不符合游戏逻辑的行为指令（如挖掘不存在的方块）。
\end{enumerate}

\subsection{未来工作}

\begin{enumerate}
    \item \textbf{多智能体协作}：支持多个 AI 之间的交流与协作。
    \item \textbf{视觉感知}：集成视觉模型，使 AI 能够"看到"游戏画面。
    \item \textbf{技能学习}：参考 Voyager\cite{wang2023voyager}，通过代码生成实现新技能的自动学习。
    \item \textbf{本地模型优化}：针对 Minecraft 领域微调小型模型，降低延迟和成本。
\end{enumerate}

% ==================== 9. 结论 ====================
\section{结论}

本文提出了 MC-agint，一种基于大语言模型的 Minecraft 自主智能体架构。通过将聊天与行为解耦、引入三层记忆系统和记忆蒸馏机制，MC-agint 实现了 AI 玩家在虚拟世界中的自然交互与自主生存。系统以 Minecraft Forge Mod 形式实现，支持 OpenAI、DeepSeek、Ollama 等兼容接口，具有良好的可扩展性。

MC-agint 的核心设计理念是\textbf{简单优先}：不追求复杂的框架，而是充分利用 LLM 本身的智能来弥补规则系统的不足。聊天时让 AI 直接说人话，记忆管理交给 LLM 蒸馏，行为规划通过简洁的 JSON 协议完成。这种设计在工程实现的简洁性和用户体验的自然性之间取得了良好的平衡。

% ==================== 参考文献 ====================
\newpage
\begin{thebibliography}{99}

\bibitem{minecraft2023}
Mojang Studios.
\textit{Minecraft Sales Surpass 300 Million}.
2023.

\bibitem{brown2020language}
Tom Brown, Benjamin Mann, Nick Ryder, et al.
\textit{Language Models are Few-Shot Learners}.
NeurIPS, 2020.

\bibitem{ouyang2022training}
Long Ouyang, Jeff Wu, Xu Jiang, et al.
\textit{Training language models to follow instructions with human feedback}.
NeurIPS, 2022.

\bibitem{millington2019artificial}
Ian Millington, John Funge.
\textit{Artificial Intelligence for Games}.
CRC Press, 2019.

\bibitem{mnih2015human}
Volodymyr Mnih, Koray Kavukcuoglu, David Silver, et al.
\textit{Human-level control through deep reinforcement learning}.
Nature, 518(7540):529--533, 2015.

\bibitem{johnson2016malmo}
Matthew Johnson, Katja Hofmann, Tim Hutton, David Bignell.
\textit{The Malmo Platform for Artificial Intelligence Experimentation}.
IJCAI, 2016.

\bibitem{park2023generative}
Joon Sung Park, Joseph C. O'Brien, Carrie J. Cai, et al.
\textit{Generative Agents: Interactive Simulacra of Human Behavior}.
UIST, 2023.

\bibitem{wang2023voyager}
Guanzhi Wang, Yuqi Xie, Yunfan Jiang, et al.
\textit{Voyager: An Open-Ended Embodied Agent with Large Language Models}.
arXiv:2305.16291, 2023.

\bibitem{autogpt2023}
Significant Gravitas.
\textit{AutoGPT}.
\url{https://github.com/Significant-Gravitas/AutoGPT}, 2023.

\bibitem{babyagi2023}
Yohei Nakajima.
\textit{BabyAGI}.
\url{https://github.com/yoheinakajima/babyagi}, 2023.

\bibitem{langchain2023}
LangChain.
\textit{LangChain Memory Documentation}.
\url{https://docs.langchain.com/oss/python/concepts/memory}, 2023.

\bibitem{coala2023}
Theodore R. Sumers, Shunyu Yao, Karthik Narasimhan, Thomas L. Griffiths.
\textit{Cognitive Architectures for Language Agents}.
arXiv:2309.02427, 2023.

\end{thebibliography}

\end{document}
